
\subsection{相对 spin 结构的改变与定向}
\label{subsec:8-1-4}

本节余下部分假设 $L$ 连通。在命题~\ref{prop:8-1-6} 中，我们定义了
$L$ 上相对 spin 结构的稳定共轭类所成集合
$\operatorname{Spin}(M,L)$。群 $H^2(M,L;\mathbb Z_2)$ 在其上单传递地
作用。我们把这一作用写成
\[
 (\mathfrak x,[V,st,\sigma])
 \longmapsto \mathfrak x\cdot[V,st,\sigma].
\]

设 $[V,st,\sigma]\in\operatorname{Spin}(M,L)$。由于
$st=w_2(V)$ 已由 $V$ 自动确定，以下简写为 $[V,\sigma]$。
设 $\beta\in\pi_2(M,L)$。模空间
$\mathcal M^{\mathrm{main}}_{k+1}(\beta)$ 带有仓西结构，而
$[V,\sigma]$ 诱导其定向。需要指明这一定向时，我们写作
\[
 \mathcal M^{\mathrm{main}}_{k+1}(\beta)^{[V,\sigma]}.
\]

注意，$\beta\in\pi_2(M,L)$ 确定一个类
$[\beta]\in H_2(M,L;\mathbb Z_2)$，故对
$\mathfrak x\in H^2(M,L;\mathbb Z_2)$，有配对
$[\beta]\cap\mathfrak x\in\mathbb Z_2$。下面的结果是
\textbf{[Cho04I]} 定理 6.4 的一个小修改。

\begin{proposition}{8.1.16}\label{prop:8-1-16}
恒等映射
\[
 \mathcal M^{\mathrm{main}}_{k+1}(\beta)^{[V,\sigma]}
 \longrightarrow
 \mathcal M^{\mathrm{main}}_{k+1}(\beta)^{
   \mathfrak x\cdot[V,\sigma]}
\]
当且仅当 $[\beta]\cap\mathfrak x=0$ 时保定向。%
\translatornote{源书在右端把上标 $\mathrm{main}$ 误印为
$\mathfrak x\cdot\mathrm{main}$；这里的恒等映射只比较同一主模空间
在 $[V,\sigma]$ 与 $\mathfrak x\cdot[V,\sigma]$ 所诱导的两个定向，
故恢复为 $\mathcal M^{\mathrm{main}}_{k+1}(\beta)$。}
\end{proposition}

\begin{proof}
只须在一点
\[
 w_0\colon(D^2,\partial D^2)\longrightarrow(M,L)
\]
比较线性化算子的指标丛定向。令
$\ell=w_0|_{\partial D^2}$。

$\mathfrak x$ 在 $M_{[3]}/L_{[2]}$ 上确定一个可定向向量丛
$E_{\mathfrak x}$。把它拉回到 $M_{[3]}$，得到向量丛
$V_{\mathfrak x}$；后者在 $L_{[2]}$ 上带有平凡化。并且
\[
 \mathfrak x\cdot[V,\sigma]
 =[V\oplus V_{\mathfrak x},\sigma\oplus\mathrm{id}]
\]
（见命题~\ref{prop:8-1-6} 的证明）。

不妨假设 $w_0(D^2)\subset M_{[3]}$，于是
$w_0^*V_{\mathfrak x}$ 有定义。由于 $D^2$ 可缩，
$w_0^*V_{\mathfrak x}$ 可平凡化，而且这一平凡化在同伦意义下唯一。
它在边界上诱导 $\ell^*V_{\mathfrak x}$ 的平凡化。

记
\[
 \pi\colon M_{[3]}\longrightarrow M_{[3]}/L_{[2]}.
\]
当且仅当 $[\beta]\cap\mathfrak x=0$ 时，上述边界平凡化与
$V_{\mathfrak x}=\pi^*E_{\mathfrak x}$ 在
$\ell^*V_{\mathfrak x}$ 上给出的平凡化相合。

另一方面，$\sigma$ 在 $\ell^*(TL\oplus V)$ 上诱导 spin 结构；
换言之，$\sigma$ 诱导 $\ell^*(TL\oplus V)$ 的一个稳定平凡化。
此外，$w_0^*V$ 平凡。合并这两个稳定平凡化，便得到
$\ell^*(TL)$ 的一个稳定平凡化。若把 $[V,\sigma]$ 换成
$[V\oplus V_{\mathfrak x},\sigma\oplus1]$，则当且仅当
$[\beta]\cap\mathfrak x\neq0$ 时，$\ell^*(TL)$ 上所诱导的
平凡化发生改变。

考虑算子
\[
 D_{w_0}\bar\partial\colon
 W^{1,p}(D^2,\partial D^2;w_0^*TM,\ell^*TL)
 \longrightarrow
 L^p\bigl(D^2;w_0^*TM\otimes\Lambda^{0,1}(D^2)\bigr).
\]
由于 $D^2$ 可缩，$w_0^*TM$ 可平凡化。如上所述，一个相对 spin
结构确定 $\ell^*TL$ 的一个平凡化同伦类。因此，指标丛的定向由
命题~\ref{prop:8-1-4} 确定，而后者确定
$\mathcal M^{\mathrm{main}}_{k+1}(\beta)$ 的定向。

下面说明：当 $\ell^*(TL)$ 的稳定平凡化改变时，该定向也改变。
沿用第~\ref{subsec:8-1-2} 小节的论证。考虑两个不同的稳定平凡化
同伦类
\[
 I_1,I_2\colon\ell^*TL\xrightarrow{\ \cong\ }S^1\times\mathbb R^n.
\]
它们的差由回路
\[
 \gamma\colon S^1=\partial D^2\longrightarrow SO(n)
\]
给出。由于复合
\[
 S^1\longrightarrow SO(n)\longrightarrow U(n)
\]
零同伦，它可延拓为 $u\colon D^2\to U(n)$。利用%
\translatornote{源书把映射对的靶误印为 $(U(n),SO(2n))$；
但边界映射 $\gamma$ 取值于 $SO(n)\subset U(n)$，且
$SO(2n)$ 并非 $U(n)$ 的子空间，故这里恢复为 $(U(n),SO(n))$。}
\[
 (u,\gamma)\colon(D^2,\partial D^2)
 \longrightarrow(U(n),SO(n)),
\]
可以得到 $S^1$ 上的一族丛对
\[
 (E,F)=[0,1]\times(w_0^*TM,\ell^*TL)/{\sim},
\]
其中等价关系 $\sim$ 由
\[
 (u,\gamma)\colon
 \{0\}\times(w_0^*TM,\ell^*TL)
 \longrightarrow
 \{1\}\times(w_0^*TM,\ell^*TL)
\]
给出。于是我们正处于命题~\ref{prop:8-1-7} 的情形；该命题说明，
稳定平凡化的改变会反转行列式线
\[
 \det(\operatorname{Coker}D_{w_0}\bar\partial)^*
 \otimes\det\operatorname{Ker}D_{w_0}\bar\partial
\]
的定向。命题得证。
\end{proof}
